\subsection{The radial coordinate}

The exact coordinate \(s_n^\star\) appears in
Theorem~\ref{thm:main-support-law}. In the range
\(d_n\log n=o(n^2)\), the spatial estimates are cleaner in the following
\(\log n\)-based coordinate:
\begin{equation}
\label{eq:s-definition}
  s_n
  =
  \frac{\gamma_nL_n-b_n}{a_n}.
\end{equation}
The corresponding critical variance defect is
\begin{equation}
\label{eq:critical-epsilon}
  \eps_{c,n}=1-\frac{L_n}{b_n},
\end{equation}
and an \(O(1)\) change in \(s_n\) changes \(\eps_n\) on the scale
\begin{equation}
\label{eq:epsilon-window}
  \Delta\eps_n
  \asymp
  \frac{a_nL_n}{b_n^2}.
\end{equation}
Lemma~\ref{lem:exact-old-edge-coordinate} shows that \(s_n\) agrees with the
exact coordinate \(s_n^\star\) whenever \(d_n\log n=o(n^2)\). Define
\[
  \xi_{n,i}
  =
  \frac{|\widetilde Z_i|^2/2-b_n}{a_n},
  \qquad
  \delta_n=\frac{a_n}{\gamma_n}.
\]
At the normalization \eqref{eq:s-definition}, the bump height is exactly
\begin{equation}
\label{eq:bump-height-xi}
  h_i
  =
  \exp\{\delta_n(\xi_{n,i}-s_n)\}.
\end{equation}
The score radius satisfies the exact identity
\begin{equation}
\label{eq:universal-score-radius}
  |q_i|^2
  =
  \frac{|\widetilde Z_i|^2}{\gamma_n}
  =2L_n+2\delta_n(\xi_{n,i}-s_n).
\end{equation}
Consequently, $|q_i|^2=2L_n+O_{\Pbb}(\delta_n)$ for edge observations with
$\xi_{n,i}=O_{\Pbb}(1)$ whenever $s_n=O(1)$.
Thus the top bumps live on a sphere of radius \(\sqrt{2L_n}\). The dimension
affects their angular geometry and the height scale \(\delta_n\), but not
their leading score radius.

When \(d_n\gg L_n\), however,
\(\delta_n\to0\), so a radial gap of \(a_n\eta\) below threshold creates
only a height gap of order \(\delta_n\eta\). The remainder of the field must
therefore be \(o_{\Pbb}(\delta_n)\), not merely \(o_{\Pbb}(1)\).

\subsection{Consequences across dimension scales}
\label{sec:regime-map}

The exact definitions \eqref{eq:b-a-definition} are preferable in the main
theorem, but their asymptotic expansions reveal how the critical defect and
window width depend on dimension. In this subsection
\(\eps_{c,n}=1-L_n/b_n\) denotes a leading-order proxy for the critical
defect. Section~\ref{sec:full-growing-transition} restores the finite-sample
correction that becomes visible when \(d_n\asymp n^2\log n\).

\medskip
\noindent\emph{Sublogarithmic dimension.}

Suppose \(k_n=o(L_n)\). Then \(b_n/L_n\to1\), \(a_n\to1\), and
\begin{equation}
\label{eq:b-sublog}
  b_n
  =
  L_n+(k_n-1)\log L_n-\log\Gamma(k_n)
  +O\left(
    \frac{k_n}{L_n}
    \left\{1+(k_n-1)\log L_n-\log\Gamma(k_n)\right\}
  \right).
\end{equation}
When in addition \(k_n\to\infty\), Stirling's formula gives the simpler
leading term
\begin{equation}
\label{eq:b-sublog-stirling}
  b_n-L_n
  \sim
  k_n\log\frac{L_n}{k_n}.
\end{equation}
Consequently,
\begin{equation}
\label{eq:epsilon-sublog}
  \eps_{c,n}
  \sim
  \frac{k_n}{L_n}\log\frac{L_n}{k_n}
  =
  \frac{d_n}{2\log n}
  \log\frac{2\log n}{d_n},
  \qquad
  \Delta\eps_n\asymp\frac1{\log n}.
\end{equation}

For fixed \(d\ge3\), \eqref{eq:b-sublog} reduces to
\[
  b_n
  =
  \log n+\left(\frac d2-1\right)\log\log n
  -\log\Gamma(d/2)+o(1),
\]
recovering the fixed-dimensional normalization used in
Section~\ref{sec:fixed-d-patches}. If
\(d_n\to\infty\), then \(b_n-L_n\to\infty\); this makes the diffuse edge
background vanish, leaving the radial maximum as the limiting obstruction.

\begin{proof}[Justification of \eqref{eq:b-sublog}]
For \(y>k_n-1\), \eqref{eq:tail-density-ratio} and
\(1\le\overline F_n(y)/f_n(y)\le\{1-(k_n-1)/y\}^{-1}\) give
\[
  \log\overline F_n(y)
  =
  -y+(k_n-1)\log y-\log\Gamma(k_n)+O(k_n/y).
\]
Standard Gamma Chernoff bounds first give \(b_n/L_n\to1\). Substitute
\(y=b_n\), write \(b_n=L_n+r_n\), and expand
\(\log b_n=\log L_n+O(r_n/L_n)\). The resulting equation first gives
\(r_n=O((k_n-1)\log L_n-\log\Gamma(k_n)+1)\), and a second substitution
gives \eqref{eq:b-sublog}. Formula \eqref{eq:b-sublog-stirling} follows
because \(L_n/k_n\to\infty\).
\end{proof}

\medskip
\noindent\emph{Logarithmic dimension.}

Suppose \(k_n/L_n\to c\in(0,\infty)\). Let \(x_c>1\) be the unique
solution of
\begin{equation}
\label{eq:xc}
  x_c-1-\log x_c=\frac1c.
\end{equation}
Stirling's formula and the Gamma saddlepoint estimate yield
\begin{equation}
\label{eq:logarithmic-regime}
  \frac{b_n}{k_n}\to x_c,
  \qquad
  a_n\to\frac{x_c}{x_c-1},
  \qquad
  \eps_{c,n}\to1-\frac1{cx_c}.
\end{equation}
The transition therefore occurs at a nonzero constant variance defect, while
its width in \(\eps_n\) remains of order \(1/\log n\).

\medskip
\noindent\emph{Superlogarithmic dimension.}

Suppose \(k_n/L_n\to\infty\). The Gamma quantile is in its
moderate-deviation range:
\begin{equation}
\label{eq:superlog-regime}
  b_n
  =
  k_n+\sqrt{2k_nL_n}\{1+o(1)\},
  \qquad
  a_n
  =
  \sqrt{\frac{k_n}{2L_n}}\{1+o(1)\}.
\end{equation}
Hence
\begin{equation}
\label{eq:epsilon-superlog}
  \eps_{c,n}
  =
  1-\frac{L_n}{k_n}\{1+o(1)\},
  \qquad
  \Delta\eps_n
  \asymp
  \frac{\sqrt{L_n}}{k_n^{3/2}},
  \qquad
  \delta_n
  \asymp
  \sqrt{\frac{L_n}{k_n}}.
\end{equation}
The fitted variance must now be much larger than the data variance: the
critical defect tends to one. The shrinking factor \(\delta_n\) in
\eqref{eq:epsilon-superlog} is also the precision required of the
single-bump approximation.

In particular, the usual proportional-dimensional regime
\(d_n/n\to\tau\in(0,\infty)\) is a special case of the superlogarithmic
regime, not the logarithmic regime \(d_n\asymp\log n\).  In that case,
\[
  1-\eps_{c,n}
  \sim\frac{2\log n}{d_n}
  \sim\frac{2\log n}{\tau n},
  \qquad
  \Delta\eps_n\asymp\frac{\sqrt{\log n}}{n^{3/2}}.
\]
The raw squared radii are of order \(n\), but after the variance rescaling
the likelihood of one empirical bump still has to overcome its factor
\(1/n\); this is why the KKT threshold remains \(\log n\).

For completeness, all three regimes follow from the same saddlepoint formula.
If \(k_n\to\infty\), put \(x_n=b_n/k_n>1\). Uniformly at the extreme
quantile,
\begin{equation}
\label{eq:gamma-saddlepoint}
  \overline F_n(k_nx_n)
  =
  \frac{1+o(1)}{(x_n-1)\sqrt{2\pi k_n}}
  \exp\{-k_n(x_n-1-\log x_n)\}.
\end{equation}
Indeed, Stirling's formula gives the density at \(k_nx_n\), while
Lemma~\ref{lem:gamma-hazard} gives the tail-to-density ratio
\(x_n/(x_n-1)\{1+o(1)\}\). Equating the logarithm of
\eqref{eq:gamma-saddlepoint} to \(-L_n\) proves
\eqref{eq:logarithmic-regime}; when \(k_n/L_n\to\infty\), expanding
\(x-1-\log x=(x-1)^2/2+O((x-1)^3)\) proves
\eqref{eq:superlog-regime}.

Finally, \eqref{eq:centering-condition} is automatic when \(k_n=O(L_n)\).
If \(k_n/L_n\to\infty\), \eqref{eq:superlog-regime} reduces it to
\begin{equation}
\label{eq:centering-superlog}
  k_nL_n=o(n).
\end{equation}
